\chapter{Setup and constants}\label{chap:setup}

\section{The comparison problem}

We formalize the following question: what is the smallest positive constant
$c$ such that every centered real random variable $X$ satisfying
\[
  \Prb(|X| > t) \le 2 e^{-t^2/2} \qquad (t \ge 0)
\]
is dominated in convex order by the centered Gaussian law of standard
deviation~$c$?

\begin{definition}[Admissible scale and sharp constant]\label{def:endpoint}
\lean{AdmissibleScale}\leanok
\lean{cStar}\leanok
A positive number $c$ is \emph{admissible} if every centered real random
variable $X$ with the two-sided sub-Gaussian tail bound above satisfies
\[
  \E[f(X)] \le \E[f(cG)]
\]
for every convex function $f$ whose expectations on both sides are finite.
The \emph{sharp constant} $c_\star$ is the infimum of all admissible scales.
\end{definition}

\begin{definition}[Explicit constants]\label{def:constants}
\lean{t0}\leanok
\lean{A}\leanok
\lean{B}\leanok
\lean{a}\leanok
\lean{p0}\leanok
\lean{z}\leanok
\lean{c0}\leanok
\lean{gaussianScale}\leanok
The explicit constant $c_0$ is built from the sub-Gaussian tail cap
\[
  s_G(t) := \min\{1, 2e^{-t^2/2}\}, \qquad
  t_0 := \sqrt{2 \log 2},
\]
the total tail mass
\[
  A := \int_0^\infty s_G(t)\,dt, \qquad B := A/2,
\]
the unique threshold $a > t_0$ solving
\[
  H(a) = a s_G(a) + \int_a^\infty s_G(t)\,dt = B,
\]
the height $p_0 := s_G(a)$, the Gaussian tail quantile $z$ satisfying
$\overline{\Phi}(z) = p_0$, and finally
\[
  c_0 := \frac{B}{\varphi(z)}.
\]
\end{definition}

\section{Comparison and sharpness statements}

The sharp theorem splits into two halves:

\begin{theorem}[Domination and sharpness]\label{thm:public-support}
\lean{gaussianDomination_c0}\leanok
\lean{gaussianSharpness_below_c0}\leanok
\begin{enumerate}
\item \textbf{Domination.} Every centered sub-Gaussian random variable
  is convex-dominated by the Gaussian law at scale~$c_0$.
\item \textbf{Sharpness.} For every $c < c_0$, domination fails:
  the canonical witness measure $\mu^\star$ and a hinge test function
  $x \mapsto (x-u)_+$ exhibit a strict inequality in the opposite direction.
\end{enumerate}
\end{theorem}

\begin{proof}\leanok
For the domination half, the proof pushes a random variable $X$ forward to its
law, verifies that the pushforward inherits the mean-zero and tail hypotheses,
and then applies the measure-level domination result from the internal
analytic argument (Theorem~\ref{thm:cx-route}).
For sharpness, the extremal witness $\mu^\star$ is constructed explicitly and
its stop-loss transform is shown to exceed the Gaussian stop-loss at the
threshold $u = cz$.
\end{proof}
